\section{linear recursive equations}
\label{sec:ricorrenza_lineare}

Given the constants $c_0, c_1, \dots, c_n$,
with $c_n \neq 0$,
we shall say that the recursive equation:
\begin{equation}\label{eq:ricorsione_lineare}
  c_n \cdot a_{k+n} + c_{n-1} \cdot a_{k+n-1} + \dots
    + c_1 \cdot a_{k+1} + c_0 \cdot a_k = 0
\end{equation}
is a \emph{linear recursive equation}%
\mymargin{linear recursive equation}\index{equation!recursive!linear} (homogeneous) of order $n$.
The polynomial $P$ with the same coefficients:
\[
    P(\lambda) = c_n \cdot \lambda^n + c_{n-1}\cdot \lambda^{n-1} + \dots + c_1
  \cdot \lambda + c_0
\]
is called the \emph{characteristic polynomial}%
\mymargin{characteristic polynomial}\index{polynomial!characteristic}
\index{equation!recursive!linear!characteristic polynomial}%
\index{polynomial!characteristic!of a linear recursive equation}%
associated with the linear equation~\eqref{eq:ricorsione_lineare}.

For example, the Fibonacci sequence $F_k$ defined by~\eqref{eq:Fibonacci}
is a solution of the second-order linear recursive equation:
\[
   a_{k+2} - a_{k+1} - a_{k} = 0
\]
whose characteristic polynomial is
\[
  P(\lambda) = \lambda^2 - \lambda - 1.
\]

\begin{theorem}[independence of exponential sequences]
\label{th:indipendenza_successioni_esponenziali}
Let $\lambda_1, \dots, \lambda_n$ be distinct complex numbers. 
Then the sequences $\vec a_j$ defined by
\[
  \vec a_j (k) = \lambda_j^k
\]
are independent vectors in the complex vector space $\CC^\NN$ of all
sequences with complex values.
\end{theorem}
%
\begin{proof}
We must show that if a linear combination of the
sequences $\vec a_j$ is zero, then all the coefficients
are zero. 
We prove this by induction on the number $n$. 
If $n=1$, the result is trivial. 
Suppose then that we know $a_1,\dots, a_{n-1}$ are independent 
and we prove that $a_1,\dots, a_n$ are also independent.

Suppose we have $c_1, \dots, c_n\in \CC$ such that
\[
  \sum_{j=1}^n c_j \vec a_j = \vec 0 \in \CC^\NN.
\]
This means that for every $k\in \NN$ we have
\begin{equation}\label{eq:47862}
  \sum_{j=1}^n c_j \lambda_j^k = 0.
\end{equation}
If we replace $k$ with $k+1$, we have
\[
  \sum_{j=1}^n c_j \lambda_j^{k+1} = 0.
\]
If we now multiply \eqref{eq:47862} by $\lambda_n$ and
subtract the previous equation, we obtain:
\[
  \sum_{j=1}^n c_j \lambda_j^k(\lambda_n - \lambda_j) = 0.   
\]
Setting $b_j = c_j (\lambda_n - \lambda_j)$, we thus obtain,
for every $k\in \NN$
\[
 \sum_{j=1}^{n-1}  b_j \lambda_j^k = 0.
\]
By the inductive hypothesis, we know that the coefficients $b_1,\dots, b_{n-1}$ 
must all be zero, which means, 
since $\lambda_n-\lambda_j\neq 0$, that $c_1,\dots, c_{n-1}$ 
are all zero. 
But then \eqref{eq:47862} becomes 
\[
  c_n \lambda_n^k = 0  
\]
which for $k=0$ implies that $c_n=0$ as well.
\end{proof}

\begin{proof}[alternative proof]
Evaluating equation~\eqref{eq:47862} for $k=0,\dots, n-1$, we obtain
\begin{align*}
  c_1 + c_2 + \dots + c_n & = 0 \\
  c_1 \lambda_1 + c_2 \lambda_2 + \dots + c_n \lambda_n & = 0 \\
  &\quad \vdots\\
  c_1 \lambda_1^{n-1} + c_2 \lambda_2^{n-1} + \dots + c_n \lambda_n^{n-1} & = 0
\end{align*}
which is a linear system whose associated matrix is
the \emph{Vandermonde} matrix%
\mymargin{Vandermonde}\index{Vandermonde}
\index{matrix!Vandermonde}%
\[
V =
\begin{pmatrix}
  1 & 1 &  \dots & 1 \\
  \lambda_1 & \lambda_2 & \dots & \lambda_n \\
  \lambda_1^2 & \lambda_2^2 & \dots & \lambda_n^2 \\
  \vdots & \vdots & \vdots & \vdots \\
  \lambda_1^{n-1} & \lambda_2^{n-1} & \dots & \lambda_n^{n-1}
\end{pmatrix}
\]
which is known (it can be proven by induction) to have determinant
\[
  \det V = \prod_{j=2}^{n} \prod_{k=1}^{j-1}(\lambda_j-\lambda_k)
\]
and thus
is invertible if and only if
the $\lambda_j$ are distinct from each other.

Thus $c_1 = c_2 = \dots = c_n = 0$ is the only solution
of the linear system.
\end{proof}

\begin{theorem}[dimension of solutions]
\label{th:dimensione_ricorsione_lineare}
The set $W$ of all sequence solutions of equation~\eqref{eq:ricorsione_lineare}
of order $n$
is a vector subspace of dimension $n$
of the space $\CC^\NN$ of all sequences.
\end{theorem}
%
\begin{proof}
Let $\vec a \in \CC^\NN$ denote the sequence $a_k = \vec a(k)$.
The set $W$ of solutions of equation~\eqref{eq:ricorsione_lineare}
is a vector subspace of $\CC^\NN$ since if $\vec a$
and $\vec b$ are solutions, then $\vec a + \vec b$ is also a solution, and if
$\vec a$
is a solution, then $t \vec a$ is a solution for every $t\in \CC$
(this can be verified).

Consider then the operator (called \emph{jet}%
\mymargin{jet}\index{jet}) $J\colon W \to \CC^n$
defined by
\[
  J(\vec a) = \enclose{\vec a(0), \vec a(1), \dots, \vec a(n-1)}.
\]
It is easy to verify that $J$ is a linear operator since the evaluation
of a function at a point is a linear operation.

We now want to determine
\[
  \ker J = \ENCLOSE{\vec a \in W \colon J(\vec a)=0}.
\]
If $\vec a\in \ker J$, it means that the first $n$ terms
of the sequence
$a_0, a_1, \dots, a_{n-1}$ are all zero.
But if $\vec a\in W$, we know that~\eqref{eq:ricorsione_lineare} holds, and thus
$a_n$ is also zero. Applying~\eqref{eq:ricorsione_lineare} recursively,
we obtain that $a_{n+k}=0$ for every $k\in \NN$, and thus $\vec a = 0$.
Therefore, $\ker J = \ENCLOSE{\vec 0}$,
and by the known properties of linear operators, we can
assert that $J\colon W \to \CC^n$ is injective.
Thus $\dim W \le \dim \CC^n = n$.

But we can also observe that $J$ is surjective since
given the values of $a_0,a_1, \dots, a_{n-1}$, the equation
~\eqref{eq:ricorsione_lineare} allows us to determine,
inductively, all subsequent terms. Thus
given $\vec y = (y_0, \dots, y_{n-1}) \in \CC^n$,
it is possible to find $\vec a \in W$ such that $J(\vec a) = \vec y$.

Since $J\colon W\to\CC^n$ is injective and surjective,
we conclude that $\dim W = \dim \CC^n = n$.
\end{proof}

\begin{theorem}[characterization of solutions]
If the characteristic polynomial of equation~\eqref{eq:ricorsione_lineare}
has $n$ distinct zeros $\lambda_1, \dots, \lambda_n$,
then
every solution of the equation can be expressed in the form
\[
  a_k = c_1 \lambda_1^k + c_2 \lambda_2^k + \dots + c_n \lambda_n^k
\]
where $c_1, \dots, c_n$ are arbitrary constants.
\end{theorem}
%
\begin{proof}
Let $V=\CC^\NN$ be the vector space of all sequences.
We can define the \emph{shift} operator $\sigma\colon V\to V$
\[
  (\sigma \vec a)(k) = \vec a(k+1).
\]
It follows that $a_{k+n} = (\sigma^n \vec a)(k)$ where
$\sigma^n = \sigma \circ \sigma \circ \dots \circ \sigma$
is the composition of $\sigma$ with itself $n$ times.
We can thus rewrite equation~\eqref{eq:ricorsione_lineare} in the form
\[
  c_n \sigma^n \vec a + \dots + c_0 \sigma^0 \vec a = 0
\]
or
\[
  P(\sigma)\vec a = 0
\]
where $P(\sigma)$ denotes the operator $P(\sigma)\colon V\to V$
obtained by applying the characteristic polynomial $P$ to the operator $\sigma$.

If $\lambda_1,\dots,\lambda_n$ are the $n$ roots of the polynomial $P$, we can
decompose the polynomial $P$ using Ruffini's theorem~\ref{th:Ruffini}
\begin{equation}\label{eq:375628}
  P(\lambda)
  = a_n (\lambda -\lambda_1)(\lambda- \lambda_2)\cdots (\lambda - \lambda_n).
\end{equation}
The same decomposition can be done for $P(\sigma)$,
and thus equation~\eqref{eq:ricorsione_lineare}
can be written in the form:
\begin{equation}\label{eq:437343}
(\sigma - \lambda_1 I)\circ (\sigma-\lambda_2 I)\circ \dots \circ (\sigma
-\lambda_n I) \vec a = 0
\end{equation}
where $I\colon V\to V$ is the identity operator $I\vec a = \vec a$.
Since in the decomposition~\eqref{eq:375628}
the order in which the factors are multiplied does not matter,
in equation~\eqref{eq:437343}, we can apply the operators
$(\sigma -\lambda_j I)$ in any order (the operators commute).
Thus it is immediate to notice that if $\vec a$ solves
\[
  (\sigma - \lambda_j I) \vec a = 0
\]
for any $j=1,\dots, n$, then $\vec a\in W$.
But this last equation is easy to solve since
setting $a_k = \vec a(k)$, it is explicit in the form:
\[
  a_{k+1} - \lambda_j a_k = 0
\]
which is satisfied if
one chooses $a_k = \lambda_j^k$.

Thus the sequence $\vec a_j$ defined by
\[
  \vec a_j(k) = \lambda_j^k
\]
is an element of $W$.
For $j=1,\dots, n$, we have found $n$ distinct solutions
$\vec a_1, \dots, \vec a_n$. These solutions are
independent by theorem~\ref{th:indipendenza_successioni_esponenziali},
and thus, since $\dim W = n$ by theorem~\ref{th:dimensione_ricorsione_lineare},
we can assert that these solutions form a basis of $W$,
which is exactly what we wanted to prove.
\end{proof}

\begin{example}[explicit formula for the Fibonacci sequence]
\index{sequence!Fibonacci}%
\index{Fibonacci!explicit formula}%
The Fibonacci sequence $F_k$ defined by~\eqref{eq:Fibonacci}
satisfies a second-order linear recursive equation with associated polynomial
$P(\lambda) = \lambda^2 - \lambda - 1$. Solving the equation $P(\lambda)=0$,
we find the two distinct roots:
\[
  \lambda_{1,2} = \frac{1 \pm \sqrt{5}}{2}.
\]
The larger root, $\lambda_2$, corresponds to the
\emph{golden ratio}%
\mymargin{golden ratio}\index{ratio!golden} which is usually denoted by $\phi$.
\index{$\phi$}%
Since $\lambda_1 \lambda_2 = -1$, we find $\lambda_1 = -1/\phi$.
Thus, by the previous theorem, we know that there exist two constants $c_1$ and
$c_2$ such that
\[
  F_k = c_1 \frac{(-1)^k}{\phi^k} + c_2 \phi^k .
\]
Imposing the conditions $F_0=0$, $F_1=1$ (imposing $F_2 = 1$ is equivalent to
$F_0=0$),
we obtain
\[
  c_1 + c_2 = 0, \qquad - \frac{c_1}{\phi} + c_2 \phi  = 1
\]
from which it is possible to find $c_1 = - c_2 = 1/\sqrt{5}$ and thus obtain
Binet's formula for the $k$-th Fibonacci number:
\mymargin{Binet's Formula}%
\index{Fibonacci}%
\index{sequence!Fibonacci}%
\index{formula!Binet's}%
\index{Binet!formula of}%
\[
  F_k = \frac{\phi^k}{\sqrt 5} + \frac{(-1)^{k+1}}{\sqrt 5\phi^k}.
\]
It is quite surprising to observe that the values $F_k$ obtained
with this expression are all integers.
\end{example}